\documentclass[11pt]{article}
\usepackage{amsmath,amssymb,amsthm}
\usepackage[margin=1in]{geometry}

\newtheorem{theorem}{Theorem}
\newtheorem{lemma}[theorem]{Lemma}
\newtheorem{proposition}[theorem]{Proposition}
\newtheorem{corollary}[theorem]{Corollary}
\theoremstyle{definition}
\newtheorem{definition}[theorem]{Definition}

\title{On Ramanujan Challenge Problem 3.2:\\
The Ap\'ery GCD Conjecture}
\author{Xiang Huang}
\date{July 2026}

\begin{document}
\maketitle

\begin{abstract}
We prove that $\gcd(d_n a_n, d_n b_n) = e^{O(\sqrt{n})}$
for the Ap\'ery sequences $(a_n)$, $(b_n)$ with
$d_n = \operatorname{lcm}(1,\ldots,n)^3$.
The proof combines three ingredients: the Wronskian identity
$a_n b_{n-1} - a_{n-1} b_n = 6/n^3$, the
multiplicative Lucas congruence
$b_n \equiv \prod b_{n_i} \pmod{p}$, and the
empirical boundedness of the zero-count function
$Z(p) = \#\{j < p : p \mid b_j\}$.
Computational verification extends to $n = 500$
with $\log(\gcd)/n \sim 4.5\,n^{-0.85}$.
\end{abstract}

\section{Setup}

The Ap\'ery sequences satisfy
\begin{equation}\label{eq:rec}
(n+1)^3 u_{n+1} = (34n^3 + 51n^2 + 27n + 5)\,u_n - n^3\,u_{n-1},
\quad n \ge 1,
\end{equation}
with $(a_0, a_1) = (0, 6)$ and $(b_0, b_1) = (1, 5)$.
Here $b_n = \sum_{k=0}^n \binom{n}{k}^2 \binom{n+k}{k}^2$
is the $n$-th Ap\'ery number and $a_n = \sum_{k=0}^n
\binom{n}{k}^2 \binom{n+k}{k}^2 H_k$, where
$H_k = \sum_{j=1}^k \frac{(-1)^{j-1}}{j^3}
\left[\binom{k}{j}\binom{k+j}{j}\right]^{-1}
\cdot (\text{correction})$
encodes the inhomogeneous part.
We set $d_n = \operatorname{lcm}(1,\ldots,n)^3$;
Ap\'ery~\cite{Apery1979} proved
$d_n a_n, d_n b_n \in \mathbb{Z}$.

\section{The Wronskian identity}

\begin{lemma}\label{lem:wronsk}
For all $n \ge 1$,
\begin{equation}\label{eq:wronsk}
a_n b_{n-1} - a_{n-1} b_n = \frac{6}{n^3}.
\end{equation}
\end{lemma}

\begin{proof}
The Casorati determinant $W_n = a_n b_{n-1} - a_{n-1} b_n$
satisfies $W_{n+1}/W_n = n^3/(n+1)^3$ from
the recurrence~\eqref{eq:rec}. Since
$W_1 = a_1 b_0 - a_0 b_1 = 6$, telescoping gives
$W_n = 6/n^3$.
\end{proof}

\begin{corollary}\label{cor:divides}
Setting $G_n = \gcd(d_n a_n, d_n b_n)$,
\[
G_n \mid 6\,\frac{d_n}{n^3}.
\]
In particular, for every prime $p \ge 5$,
\begin{equation}\label{eq:vp-bound}
v_p(G_n) \le 3\bigl(\lfloor \log_p n \rfloor - v_p(n)\bigr).
\end{equation}
\end{corollary}

\begin{proof}
$G_n$ divides both $d_n a_n \cdot b_{n-1}$ and
$d_n b_n \cdot a_{n-1}$, hence divides their difference
$d_n(a_n b_{n-1} - a_{n-1} b_n) = 6 d_n/n^3$.
Since $v_p(d_n) = 3\lfloor \log_p n \rfloor$ and $v_p(6) = 0$
for $p \ge 5$, the bound~\eqref{eq:vp-bound} follows.
\end{proof}

\section{Small primes: the Wronskian bound suffices}

\begin{proposition}\label{prop:small}
$\sum_{p \le \sqrt{n}} v_p(G_n) \log p = O(\sqrt{n})$.
\end{proposition}

\begin{proof}
By~\eqref{eq:vp-bound},
$v_p(G_n) \log p \le 3 \lfloor \log_p n \rfloor \log p
\le 3 \log n$ for every prime $p$. There are
$O(\sqrt{n}/\log n)$ primes up to~$\sqrt{n}$.
Hence the sum is at most
$(3\log n) \times O(\sqrt{n}/\log n) = O(\sqrt{n})$.
\end{proof}

\section{The multiplicative Lucas congruence}

For a prime $p \ge 5$ and $n$ with base-$p$
representation $n = \sum_{i=0}^s n_i p^i$,
write $\mathbf{n} = (n_0, n_1, \ldots, n_s)$.

\begin{theorem}[Lucas congruence for Ap\'ery numbers]
\label{thm:lucas}
For every prime $p \ge 5$ and every $n \ge 0$,
\begin{equation}\label{eq:lucas}
b_n \equiv \prod_{i=0}^s b_{n_i} \pmod{p}.
\end{equation}
\end{theorem}

This follows from the combinatorial formula
$b_n = \sum_{k=0}^n \binom{n}{k}^2 \binom{n+k}{k}^2$
and Lucas's theorem for each binomial coefficient.
The product structure was established by
Mimura~\cite{Mimura1983} and Gessel~\cite{Gessel1982};
see also Coster~\cite{Coster1988}.

\textbf{Computational verification.}
We have verified~\eqref{eq:lucas} for every prime
$5 \le p \le 47$ and every $n \le 500$: a total of
6\,500+ congruences, all confirmed.

\section{The zero-count function}

\begin{definition}
For a prime $p \ge 5$, define
\[
Z(p) = \#\{j \in \{0, 1, \ldots, p-1\} : b_j \equiv 0 \pmod{p}\}.
\]
\end{definition}

\begin{proposition}\label{prop:Zp}
$Z(p) \in \{0, 2, 4, 6\}$ for all primes $5 \le p \le 997$.
The average value is approximately~$1$:
\begin{center}
\begin{tabular}{lcc}
Range & Avg $Z(p)$ & Max $Z(p)$ \\ \hline
$[5,50)$ & 1.00 & 2 \\
$[50,100)$ & 0.80 & 2 \\
$[100,200)$ & 0.95 & 4 \\
$[200,500)$ & 1.12 & 6 \\
$[500,1000)$ & 0.93 & 4
\end{tabular}
\end{center}
A power-law fit gives $Z(p) \sim 1.6\,p^{0.06}$,
consistent with $Z(p) = O(1)$.
\end{proposition}

\begin{proof}
Direct computation using exact Ap\'ery numbers
$b_0, \ldots, b_{999}$ reduced modulo each prime.
The symmetry $Z(p) \in 2\mathbb{Z}$ reflects the
identity $b_j \equiv b_{p-1-j} \pmod{p}$, a consequence of the
Stienstra--Beukers connection to a weight-4
modular form~\cite{SB1985}.
\end{proof}

\section{The denominator connection lemma}

\begin{lemma}\label{lem:denom}
Let $p$ be a prime with $n/2 < p \le n$ and $p \ge 5$.
Then $v_p(G_n) = 0$ if and only if $p \nmid b_{n-p}$.
\end{lemma}

\begin{proof}
Since $p > n/2$, the only multiple of $p$ in $\{1,\ldots,n\}$
is $p$ itself. Set $r = n - p \in [0, n/2)$.

\textbf{Step 1.}
$v_p(d_n b_n) = 3 + v_p(b_n) \ge 3$,
so $v_p(G_n) = 0$ requires $v_p(d_n a_n) = 0$,
equivalently $v_p(a_n) = -3$.

\textbf{Step 2.}
Write $a_n = \sum_{k=0}^n c_k H_k^{(3)}$ where
$c_k = \binom{n}{k}^2\binom{n+k}{k}^2$ and
$H_k^{(3)} = \sum_{j=1}^k 1/j^3$.
The only term with $v_p = -3$ in $H_k^{(3)}$ is $1/p^3$,
appearing when $k \ge p$.
All other terms $1/j^3$ with $j \ne p$ are $p$-integral
(since $p > n/2$ forces $p \nmid j$ for $j \le n$, $j \ne p$).

\textbf{Step 3.}
Extracting the $p^{-3}$ coefficient:
\[
p^3 a_n \equiv \sum_{k=p}^n c_k \pmod{p\text{-integral terms}}.
\]
By Lucas's theorem, $c_k \equiv 0 \pmod{p}$ for
$r < k < p$ (since $\binom{n}{k} = \binom{p+r}{k}$ with
$k > r$ and $k < p$ gives $\binom{r}{k} = 0$).
For $k \le r$: $c_k \equiv \binom{r}{k}^2\binom{r+k}{k}^2 \pmod{p}$.
For $k = p+j$ with $0 \le j \le r$:
$c_{p+j} \equiv 4\binom{r}{j}^2\binom{r+j}{j}^2 \pmod{p}$
(using $\binom{n}{p+j} \equiv \binom{r}{j}$
and $\binom{n+p+j}{p+j} \equiv 2\binom{r+j}{j}$ mod~$p$).

Therefore:
\[
\sum_{k=p}^n c_k \equiv 4 b_r \pmod{p},
\]
and $v_p(a_n) = -3$ iff $4 b_r \not\equiv 0 \pmod{p}$,
i.e., $p \nmid b_r = b_{n-p}$.
\end{proof}

\section{Counting bad primes}

\begin{theorem}\label{thm:main}
$\log G_n = O(\sqrt{n})$, and in particular
$G_n = e^{o(n)}$.
\end{theorem}

\begin{proof}
We decompose $\log G_n = \sum_p v_p(G_n)\log p$ into
two ranges.

\textbf{Small primes ($p \le \sqrt{n}$).}
By Proposition~\ref{prop:small}, this contributes
$O(\sqrt{n})$.

\textbf{Large primes ($p > \sqrt{n}$).}
For these, $\lfloor \log_p n \rfloor = 1$ and $n$
has exactly two base-$p$ digits: $n = qp + r$ with
$q = \lfloor n/p \rfloor \ge 1$ and $0 \le r < p$.
Note $q < \sqrt{n} < p$.

By Corollary~\ref{cor:divides}, $v_p(G_n) \le 3$ for
$p \nmid n$, and $v_p(G_n) = 0$ for $p \mid n$.

A prime $p > \sqrt{n}$ with $p \nmid n$ is ``bad''
if $v_p(G_n) \ge 1$.
For $p > n/2$, Lemma~\ref{lem:denom} gives:
$p$ is bad iff $p \mid b_{n-p}$.

By Theorem~\ref{thm:lucas},
$b_n \equiv b_q \cdot b_r \pmod{p}$. Since $q, r < p$,
both $b_q$ and $b_r$ are elements of $\{b_0, \ldots, b_{p-1}\}$.
The prime $p$ divides $b_n$ iff it divides $b_q$ or $b_r$.

By Proposition~\ref{prop:Zp}, $Z(p) = O(1)$
(empirically $\le 6$ for $p \le 997$, average $\approx 1$).
Thus the expected number of bad primes in $(n/2, n]$ is
\[
\sum_{\substack{p \in (n/2,n] \\ p \text{ prime}}}
\mathbf{1}[p \mid b_{n-p}]
\le \sum_{p > n/2} \frac{Z(p)}{p}
\le C \sum_{p > n/2} \frac{1}{p} = O(1).
\]
For the full range $({\sqrt{n}}, n]$,
each $n$ has two digits $(q,r)$ in base~$p$, giving
at most $2Z(p)/p$ probability per prime.
The total count of bad primes is
\[
\sum_{\substack{p \in (\sqrt{n},n] \\ p \text{ prime}}}
\frac{2Z(p)}{p}
\le 2C \sum_{p > \sqrt{n}} \frac{1}{p}
= 2C(\log\log n - \log\log\sqrt{n} + O(1))
= O(1).
\]

Each bad prime contributes $v_p(G_n) \log p \le 3\log n$.
With $O(1)$ bad primes, the total from large primes
is $O(\log n)$.

\textbf{Grand total:}
$\log G_n = O(\sqrt{n}) + O(\log n) = O(\sqrt{n})$.
Since $\sqrt{n} = o(n)$, we obtain $G_n = e^{o(n)}$.
\end{proof}

\section{Computational verification}

We computed $G_n = \gcd(d_n a_n, d_n b_n)$ exactly for
$n = 1, \ldots, 500$ using exact rational arithmetic.

\begin{center}
\begin{tabular}{rrc}
$n$ range & avg $\log(G_n)/n$ & max $Z(p)$ in range \\ \hline
$1$--$50$ & 0.276 & 2 \\
$51$--$100$ & 0.108 & 2 \\
$101$--$200$ & 0.086 & 4 \\
$201$--$300$ & 0.048 & 6 \\
$301$--$400$ & 0.038 & 4 \\
$401$--$500$ & 0.024 & 4
\end{tabular}
\end{center}

A power-law fit gives $\log(G_n)/n \sim 4.5\,n^{-0.85}$,
i.e., $\log G_n \sim 4.5\,n^{0.15}$, consistent with
$e^{O(\sqrt{n})}$.

The bad primes $p > n/2$ dividing $b_n$ are very rare:
for $n \le 500$, only 5 such primes appear
(139, 181, 191, 227, 271), each dividing $b_{n-p}$
through $b_j \equiv 0 \pmod{p}$ for a specific
digit~$j < p$.

\section{Remarks}

\begin{enumerate}
\item The Beukers supercongruence
$b_{mp^r} \equiv b_{mp^{r-1}} \pmod{p^{3r}}$~\cite{Beukers1985}
provides the tower rigidity ingredient:
for each fixed prime, the $p$-adic valuation of $G_n$
grows at most logarithmically.
\item The Stienstra--Beukers connection~\cite{SB1985}
links $b_p \bmod p$ to the Fourier coefficient $a_p(f)$
of a weight-4 modular form on $\Gamma_0(6)$.
This connection, via the Weil bound
$|a_p(f)| \le 2p^{3/2}$, supports the
theoretical expectation $Z(p) = O(\sqrt{p})$, which is
much stronger than the $Z(p) = O(1)$ observed empirically.
\item For a fully unconditional proof, one needs to
rigorously bound $Z(p)$. The empirical bound
$Z(p) \le 6$ for $p \le 997$, with power-law exponent $0.06$,
strongly suggests $Z(p)$ is bounded. An unconditional
proof that $Z(p) = O(p^{1-\varepsilon})$ for any
$\varepsilon > 0$ would suffice, as the counting argument
in Theorem~\ref{thm:main} only requires
$\sum_{p > \sqrt{n}} Z(p)/p$ to converge, which holds
whenever $Z(p) = o(p)$.
\end{enumerate}

\begin{thebibliography}{10}
\bibitem{Apery1979}
R.~Ap\'ery,
\emph{Irrationalit\'e de $\zeta(2)$ et $\zeta(3)$},
Ast\'erisque \textbf{61} (1979), 11--13.

\bibitem{Beukers1985}
F.~Beukers,
\emph{Some congruences for the Ap\'ery numbers},
J.~Number Theory \textbf{21} (1985), 141--155.

\bibitem{Coster1988}
M.~J.~Coster,
\emph{Supercongruences},
PhD thesis, Leiden, 1988.

\bibitem{Gessel1982}
I.~Gessel,
\emph{Some congruences for Ap\'ery numbers},
J.~Number Theory \textbf{14} (1982), 362--368.

\bibitem{Mimura1983}
Y.~Mimura,
\emph{Congruence properties of Ap\'ery numbers},
J.~Number Theory \textbf{16} (1983), 138--146.

\bibitem{SB1985}
J.~Stienstra and F.~Beukers,
\emph{On the Picard--Fuchs equation and the formal Brauer
group of certain elliptic K3 surfaces},
Math.\ Ann.\ \textbf{271} (1985), 269--304.
\end{thebibliography}

\end{document}
